\documentclass{situationist-preprint}

\title{The Paradox of Hedonism}
\msid{EXI-07-b}
\msfamily{EXI-07}
\msdiscipline{Meaning, Death and the Absurd}
\mstype{Companion manuscript}
\mskeywords{paradox of hedonism; happiness; indirect self-defeat; sophisticated hedonism; final ends}

\begin{document}
\maketitle

\begin{abstract}
The paradox of hedonism says that the pursuit of happiness defeats itself. It is often dismissed as either trivial (monitoring one's mood is distracting) or false (a sophisticated hedonist can adopt other ends as means). We argue that, in the form given by EXI-07-a, it is neither. It is a case of indirect self-defeat in Parfit's sense, and it rests on a structural feature of final ends rather than on a psychological accident. We show that the sophisticated hedonist faces a dilemma: either she values her other ends for their own sake, and then happiness is not her sole final end, or she does not, and then Lemma 3.2 of EXI-07-a applies to her. Recent empirical work is consistent with the result but is not needed for it.
\end{abstract}

\section{The paradox as indirect self-defeat}

\citet[II.iii]{Sidgwick1907} observed that the impulse towards pleasure, if too predominant, defeats its own aim, and \citet[ch.~V]{Mill1873} reported that those are happy who have their minds fixed on some object other than their own happiness. Two readings must be distinguished. On the \emph{attentional} reading, attending to one's happiness while engaged in an activity disrupts the absorption on which the enjoyment depends. On the \emph{structural} reading, the trouble lies not in attention but in what the agent values: an agent whose sole final end is her own happiness (Definition~2.2 of EXI-07-a) lacks the object-directed caring on which engagement depends (Principle~P2 of EXI-07-a).

The attentional reading is true but weak: it is answered by attending less. The structural reading is what Theorem~3.3 of EXI-07-a establishes, and it is a case of what \citet[\S1]{Parfit1984} calls \emph{indirect} self-defeat. A theory or policy is indirectly self-defeating if, when an agent successfully follows it, its aims are worse achieved than if she had followed another. The agent who aims directly at happiness does not fail by mistake or weakness; she fails because she succeeds in making happiness her sole final end.

\begin{proposition}\label{p:indirect}
The policy of aiming directly at happiness is indirectly self-defeating: its successful execution yields less happiness than the policy of having final ends other than one's happiness.
\end{proposition}

\begin{proof}
Successful execution of the policy is aiming directly at happiness in the sense of Definition~2.2 of EXI-07-a. By Theorem~3.3 of EXI-07-a, an agent who does so attains less happiness than she would if she had other final ends. The shortfall is not due to a failure to execute the policy, since it follows from its successful execution by Lemma~3.2 of EXI-07-a.
\end{proof}

\section{The sophisticated hedonist}

\citet{Railton1984} describes a sophisticated hedonist who, knowing the paradox, cultivates commitments to friends, work and causes, while holding that they are justified by their contribution to happiness. Does she escape?

\begin{proposition}\label{p:dilemma}
Either the sophisticated hedonist values the objects of her commitments for their own sake, in which case happiness is not her sole final end; or she does not, in which case she forgoes the happiness that depends on engagement.
\end{proposition}

\begin{proof}
If she values her friends' good for its own sake, she desires something other than happiness not only as a means to it, and by Definition~2.1 of EXI-07-a her happiness is not her sole final end. If she does not, then by Principle~P2 of EXI-07-a she lacks the caring that the forms of engagement in Principle~P1 require, and Theorem~3.3 of EXI-07-a applies.
\end{proof}

The first horn is not a defeat for the sophisticated hedonist as Railton describes her; it is a concession to the conclusion of EXI-07. She may hold the \emph{theory} that happiness is what ultimately justifies her commitments, but she cannot hold it as her \emph{aim}, the end that her practical reasoning serves, without passing to the second horn. Happiness then figures in her life as the by-product of Definition~2.3 of EXI-07-a, not as its goal, which is what Corollary~3.4 of EXI-07-a asserts.

\begin{objection}[Empirical contingency]
Whether pursuing happiness reduces it is an empirical question. Some studies suggest that valuing happiness highly is associated with lower well-being \citep{Mauss2011}, but such findings are correlational and culturally variable.
\end{objection}

\begin{reply}
Proposition~\ref{p:indirect} does not rest on these findings. It rests on Principles~P1 and P2 of EXI-07-a, which concern what engagement is: one cannot be absorbed in a game whose outcome does not matter to one. The findings are consistent with the result and illustrate the attentional reading, but the structural reading would survive their failure to replicate.
\end{reply}

\section{Conclusion}

The paradox of hedonism is neither a curiosity of attention nor a refutable empirical generalization. It is the indirect self-defeat of a policy about final ends, and the sophisticated hedonist escapes it only by ceasing to treat happiness as the goal of life.

\begin{thebibliography}{9}\small
\bibitem[Mauss et~al.(2011)]{Mauss2011} Mauss, I.~B., Tamir, M., Anderson, C.~L. and Savino, N.~S. (2011). Can seeking happiness make people unhappy? Paradoxical effects of valuing happiness. \emph{Emotion}, 11(4), 807--815.
\bibitem[Mill(1873)]{Mill1873} Mill, J.~S. (1873). \emph{Autobiography}. Longmans, Green, Reader and Dyer.
\bibitem[Parfit(1984)]{Parfit1984} Parfit, D. (1984). \emph{Reasons and Persons}. Clarendon Press.
\bibitem[Railton(1984)]{Railton1984} Railton, P. (1984). Alienation, consequentialism, and the demands of morality. \emph{Philosophy \& Public Affairs}, 13(2), 134--171.
\bibitem[Sidgwick(1907)]{Sidgwick1907} Sidgwick, H. (1907). \emph{The Methods of Ethics}, 7th ed. Macmillan.
\end{thebibliography}

\end{document}
